\documentclass
[kulak] % [kulak|kul|kuleuven-brugge,handout|outline|]
{kulakbeamer}

\usepackage[dutch]{babel}
\usepackage[T1]{fontenc}

\title{Werken met Git en Github vanuit VSCode}
\author{Stijn Rebry}
\date{Academiejaar 2026 -- 2027}

\definecolor{GitOrange}{RGB}{240,80,50}
\definecolor{GitHubGray}{RGB}{36,41,46}
\definecolor{VSCodeBlue}{RGB}{0,122,204}

\renewcommand{\outlineTitle}{Overzicht} % enkel met optie outline

\begin{document}
	
	\begin{titleframe}
		\titlepage
	\end{titleframe}

    \begin{frame}{Inleiding}
		\begin{itemize}
            \item \textbf{Versiebeheer}
            \begin{itemize}
                \item Informatie bewaren en herstellen
                \item Samenwerken aan hetzelfde project
                \item Gestructureerd werken als team
            \end{itemize} \pause
			\item \textbf{Git} bundelt in een repository
            \begin{itemize}
                \item projectbestanden
                \item samen met geschiedenis van aanpassingen en
                \item parallelle versies van het project
            \end{itemize} 
			\item \textbf{GitHub} bewaart repositories online
            \begin{itemize}
                \item backup van actuele, vorige en parallelle versies
                \item samenwerking
            \end{itemize}
			\item \textbf{VSCode} is de grafische interface van waaruit we Git aansturen.
		\end{itemize}
	\end{frame}

	\begin{outlineframe}[Overzicht]
		\tableofcontents
	\end{outlineframe}
	
	% ======================================================================
	\section{Local repository, commits}
	% ======================================================================
	
	\begin{frame}{Lokale map maken en openen in VSCode}
		\begin{columns}[T]
			\column{0.48\textwidth}
			\begin{block}{1. Windows Explorer}
			\begin{itemize}
				\item Maak map \texttt{C:\textbackslash Projecten}
				\item Maak submap \texttt{gittest}
			\end{itemize}
			\end{block}
			\begin{block}{2. Visual Studio Code}
			\begin{itemize}
				\item Open map \texttt{C:\textbackslash Projecten\textbackslash gittest}
				\item Maak nieuw bestand \texttt{gegevens.txt}\\
				% \begin{verbatim}
                \texttt{Naam: Stijn Rebry}\\
                \texttt{Opleiding: Wiskunde}
                % \end{verbatim}
				\item Sla bestand op
			\end{itemize}
            \end{block}
			
			\column{0.48\textwidth}
			\begin{alertblock}{Gebruik lokale map, geen cloudopslag.}
			\begin{itemize}
				\item Versieconflicten
				\item Corrupte bestanden
				\item Prestatieproblemen
			\end{itemize}
		\end{alertblock}

		\begin{exampleblock}{Nuttige sneltoetsen}
			\begin{itemize}
				\item New file: \texttt{ctrl+N}
				\item Save file: \texttt{ctrl+S}
			\end{itemize}
		\end{exampleblock}
		\end{columns}
	\end{frame}

	\begin{frame}{Lokaal repository initialiseren en eerste commit}
		\begin{columns}[T]
		
			\column{0.48\textwidth}
			\begin{block}{3. Source control}
				\begin{itemize}
					\item Open \texttt{Source Control}
					\item Klik \texttt{Trust}
					\item Klik \texttt{Initialize Repository}
				\end{itemize}
			\end{block}

			\begin{block}<2>{4. Source Control, Changes}
				\begin{itemize}
					\item Schrijf een commit message
					\item Klik \texttt{Commit}
					\item Kies \texttt{Always}
				\end{itemize}
			\end{block}
	
			\column{0.48\textwidth}
			\begin{exampleblock}{Changes}
				\begin{itemize}
					\item Lijst van aanpassingen:
					\item \texttt{gegevens.txt A}
					\item[]
				\end{itemize}
			\end{exampleblock}

			\begin{exampleblock}<2>{Graph}
			\begin{itemize}
				\item Lijst van commits
				\item Repository bevat nu één versie
			\end{itemize}
		\end{exampleblock}
		\end{columns}
	\end{frame}
	
\begin{frame}{}
		\begin{columns}[T]
			\column{0.48\textwidth}
			\begin{block}{5. \texttt{gegevens.txt}}
				\begin{itemize}
				\item Voeg woonplaats toe, sla op
				\item Bekijk \texttt{Changes} en \texttt{Graph}
				\item Maak een nieuwe commit
				\end{itemize}
			\end{block}

			\begin{block}<2>{6. Bestanden toevoegen en verwijderen}
				\begin{itemize}
				\item Maak \texttt{brol.txt}
				\item Maak \texttt{bestanden.txt}
				\item Commit
				\item Verwijder \texttt{brol.txt}
				\item Pas \texttt{bestanden.txt} aan
				\item Commit opnieuw
				\end{itemize}
			\end{block}
	
			\column{0.48\textwidth}
			\begin{exampleblock}<2>{Graph}
				Elke commit toont veranderingen:
				\begin{itemize}
					\item \texttt{M} -- Modified file
					\item \texttt{A} -- Added file
					\item \texttt{D} -- Deleted file
				\end{itemize}
			\end{exampleblock}

			\begin{exampleblock}<2>{Editor}
			Klik op Modified file:
				\begin{itemize}
				\item Toevoegingen (groen)
				\item Verwijderde tekst (rood)
				\end{itemize}
		\end{exampleblock}
		\end{columns}
	\end{frame}

\end{document}

\begin{frame}{}
		\begin{columns}[T]
			\column{0.48\textwidth}
			\begin{block}{}
				\begin{itemize}
					\item 
					\item 
					\item 
				\end{itemize}
			\end{block}

			\begin{block}<2>{}
				\begin{itemize}
					\item 
					\item 
					\item 
				\end{itemize}
			\end{block}
	
			\column{0.48\textwidth}
			\begin{exampleblock}{}
				\begin{itemize}
					\item
					\item
				\end{itemize}
			\end{exampleblock}

			\begin{exampleblock}<2>{}
			\begin{itemize}
				\item
				\item
				\item
			\end{itemize}
		\end{exampleblock}
		\end{columns}
	\end{frame}

	
	\begin{frame}{Terminologie}
		\begin{columns}[T]
			\column{0.48\textwidth}
			\textbf{Nieuwe begrippen}
			
			\begin{itemize}
				\item Werkmap
				\item Commit
				\item Repository
			\end{itemize}
			
			\column{0.48\textwidth}
			\textbf{Betekenis}
			
			\begin{itemize}
				\item Werkmap = huidige bestanden
				\item Commit = momentopname
				\item Repository =
				verzameling van alle commits
			\end{itemize}
		\end{columns}
	\end{frame}
	
	% ======================================================================
	\section{Repository publiceren op GitHub}
	% ======================================================================
	

	
\begin{frame}{}
		\begin{columns}[T]
			\column{0.48\textwidth}
			\begin{block}{}
				\begin{itemize}
					\item 
					\item 
					\item 
				\end{itemize}
			\end{block}

			\begin{block}<2>{}
				\begin{itemize}
					\item 
					\item 
					\item 
				\end{itemize}
			\end{block}
	
			\column{0.48\textwidth}
			\begin{exampleblock}{}
				\begin{itemize}
					\item
					\item
				\end{itemize}
			\end{exampleblock}

			\begin{exampleblock}<2>{}
			\begin{itemize}
				\item
				\item
				\item
			\end{itemize}
		\end{exampleblock}
		\end{columns}
	\end{frame}

\end{document}

	\begin{frame}{GitHub-account aanmaken}
		\begin{columns}[T]
			\column{0.48\textwidth}
			\textbf{Studenten}
			
			\begin{itemize}
				\item Surf naar \texttt{github.com}
				\item Klik \texttt{Sign Up}
				\item Maak een account aan
				\item Meld aan
			\end{itemize}
			
			\column{0.48\textwidth}
			\textbf{Opzet en resultaat}
			
			GitHub zal de online versie van de repository bewaren.
			
			\begin{itemize}
				\item Voorbereiding op back-up
				\item Voorbereiding op samenwerking
				\item Toegang tot online repositories
			\end{itemize}
		\end{columns}
	\end{frame}
	
	\begin{frame}{Repository publiceren vanuit VSCode}
		\begin{columns}[T]
			\column{0.48\textwidth}
			\textbf{Studenten}
			
			\begin{itemize}
				\item Klik \texttt{Publish Branch}
				\item Kies \texttt{Private Repository}
				\item Meld aan indien nodig
				\item Bevestig
			\end{itemize}
			
			\column{0.48\textwidth}
			\textbf{Opzet en resultaat}
			
			De lokale repository krijgt een online kopie.
			
			\begin{itemize}
				\item GitHub maakt een repository aan
				\item Alle bestaande commits worden gepubliceerd
				\item Lokale en online versie zijn identiek
			\end{itemize}
		\end{columns}
	\end{frame}
	
	\begin{frame}{Repository bekijken op GitHub}
		\begin{columns}[T]
			\column{0.48\textwidth}
			\textbf{Studenten}
			
			\begin{itemize}
				\item Open \texttt{Repositories}
				\item Open \texttt{gittest}
				\item Controleer de bestanden
			\end{itemize}
			
			\column{0.48\textwidth}
			\textbf{Opzet en resultaat}
			
			Verifiëren dat de publicatie geslaagd is.
			
			\begin{itemize}
				\item \texttt{gegevens.txt} zichtbaar
				\item \texttt{bestanden.txt} zichtbaar
				\item Zelfde inhoud als lokaal
			\end{itemize}
		\end{columns}
	\end{frame}
	
	\begin{frame}{Nieuwe wijzigingen lokaal toevoegen}
		\begin{columns}[T]
			\column{0.48\textwidth}
			\textbf{Studenten}
			
			\begin{itemize}
				\item Voeg favoriete kleur toe
				\item Sla op
				\item Maak een commit
				\item Vergelijk lokaal en GitHub
			\end{itemize}
			
			\column{0.48\textwidth}
			\textbf{Opzet en resultaat}
			
			Committen is niet hetzelfde als publiceren.
			
			\begin{itemize}
				\item Lokale bestanden zijn aangepast
				\item Nieuwe commit zichtbaar in Graph
				\item Wijziging ontbreekt nog op GitHub
			\end{itemize}
		\end{columns}
	\end{frame}
	
	\begin{frame}{Push uitvoeren}
		\begin{columns}[T]
			\column{0.48\textwidth}
			\textbf{Studenten}
			
			\begin{itemize}
				\item Klik \texttt{Sync Changes}
				\item Vernieuw GitHub
			\end{itemize}
			
			\column{0.48\textwidth}
			\textbf{Opzet en resultaat}
			
			Een push stuurt lokale commits naar GitHub.
			
			\begin{itemize}
				\item Favoriete kleur verschijnt online
				\item Local repository en GitHub zijn opnieuw gelijk
			\end{itemize}
		\end{columns}
	\end{frame}
	
	\begin{frame}{Wijzigingen rechtstreeks op GitHub aanbrengen}
		\begin{columns}[T]
			\column{0.48\textwidth}
			\textbf{Studenten}
			
			\begin{itemize}
				\item Open \texttt{gegevens.txt}
				\item Klik op het potlood-icoon
				\item Voeg hobby toe
				\item Commit de wijziging
			\end{itemize}
			
			\column{0.48\textwidth}
			\textbf{Opzet en resultaat}
			
			GitHub kan zelf ook nieuwe commits maken.
			
			\begin{itemize}
				\item Hobby verschijnt online
				\item Hobby ontbreekt lokaal
				\item GitHub bevat nu een recentere versie
			\end{itemize}
		\end{columns}
	\end{frame}
	
	\begin{frame}{Pull uitvoeren}
		\begin{columns}[T]
			\column{0.48\textwidth}
			\textbf{Studenten}
			
			\begin{itemize}
				\item Klik \texttt{Sync Changes}
				\item Controleer het bestand lokaal
			\end{itemize}
			
			\column{0.48\textwidth}
			\textbf{Opzet en resultaat}
			
			Een pull haalt commits van GitHub op.
			
			\begin{itemize}
				\item Hobby verschijnt lokaal
				\item Laatste online commit wordt opgehaald
				\item Beide repositories zijn opnieuw identiek
			\end{itemize}
		\end{columns}
	\end{frame}
	
	\begin{frame}{Terminologie}
		\begin{columns}[T]
			\column{0.48\textwidth}
			\textbf{Nieuwe begrippen}
			
			\begin{itemize}
				\item Local repository
				\item Remote repository
				\item Push
				\item Pull
			\end{itemize}
			
			\column{0.48\textwidth}
			\textbf{Betekenis}
			
			\begin{itemize}
				\item Lokaal = eigen computer
				\item Remote = GitHub
				\item Push = lokaal \(\rightarrow\) GitHub
				\item Pull = GitHub \(\rightarrow\) lokaal
			\end{itemize}
		\end{columns}
	\end{frame}	
	
	\section{Werken met branches}
	
\begin{frame}{Werken met branches}
	
	\begin{columns}[T]
		
		\column{0.48\textwidth}
		\textcolor{VSCodeBlue}{\textbf{Situatie}}
		
		Tot nu toe:
		
		\begin{itemize}
			\item één repository
			\item één branch
			\item alle commits op \texttt{main}
		\end{itemize}
		
		\column{0.48\textwidth}
		\textcolor{GitOrange}{\textbf{Waarom branches?}}
		
		\begin{itemize}
			\item wijzigingen uitproberen
			\item experimenteren
			\item samenwerken
			\item stabiele versie bewaren
		\end{itemize}
		
		\medskip
		
		\textcolor{GitHubGray}{\textbf{Belangrijk}}
		
		Een branch is een alternatieve versie van dezelfde repository.
		
	\end{columns}
	
\end{frame}


\begin{frame}{De huidige branch bekijken}
	
	\begin{columns}[T]
		
		\column{0.48\textwidth}
		
		\textcolor{VSCodeBlue}{\textbf{Wat doen we?}}
		
		\begin{itemize}
			\item Kijk linksonder in VSCode.
			\item Zoek de branchnaam.
			\item Controleer dat deze \texttt{main} is.
			\item Verifieer dezelfde branch op GitHub.
		\end{itemize}
		
		\column{0.48\textwidth}
		
		\textcolor{GitOrange}{\textbf{Waarom?}}
		
		Elke commit behoort tot een branch.
		
		\medskip
		
		\textcolor{GitHubGray}{\textbf{Resultaat}}
		
		\begin{itemize}
			\item actieve branch = \texttt{main}
			\item lokale en remote branch zijn identiek
			\item alle commits staan momenteel op \texttt{main}
		\end{itemize}
		
	\end{columns}
	
\end{frame}

\begin{frame}{Een nieuwe branch maken}
	
	\begin{columns}[T]
		
		\column{0.48\textwidth}
		
		\textcolor{VSCodeBlue}{\textbf{Wat doen we?}}
		
		\begin{itemize}
			\item Klik op \texttt{main}.
			\item Kies \texttt{Create New Branch}.
			\item Geef naam \texttt{vakken}.
			\item Wissel tussen \texttt{main} en \texttt{vakken}.
		\end{itemize}
		
		\column{0.48\textwidth}
		
		\textcolor{GitOrange}{\textbf{Waarom?}}
		
		Een branch vertrekt vanuit een bestaande versie van het project.
		
		\medskip
		
		\textcolor{GitHubGray}{\textbf{Resultaat}}
		
		\begin{itemize}
			\item branch \texttt{vakken} bestaat lokaal
			\item branch bevat dezelfde bestanden als \texttt{main}
			\item nog niet zichtbaar op GitHub
		\end{itemize}
		
	\end{columns}
	
\end{frame}


\begin{frame}{Wijziging maken op een branch}
	
	\begin{columns}[T]
		
		\column{0.48\textwidth}
		
		\textcolor{VSCodeBlue}{\textbf{Wat doen we?}}
		
		\begin{itemize}
			\item Selecteer branch \texttt{woonplaatsen}.
			\item Maak \texttt{woonplaatsen.txt}.
			\item Voeg inhoud toe.
			\item Commit.
			\item Toon alle branches in de Graph.
			\item Wissel tussen branches.
		\end{itemize}
		
		\column{0.48\textwidth}
		
		\textcolor{GitOrange}{\textbf{Waarom?}}
		
		Branches kunnen onafhankelijk evolueren.
		
		\medskip
		
		\textcolor{GitHubGray}{\textbf{Resultaat}}
		
		\begin{itemize}
			\item nieuwe file enkel zichtbaar in \texttt{woonplaatsen}
			\item Graph splitst op in meerdere kleuren
			\item elke branch heeft eigen commits
			\item andere branch blijft ongewijzigd
		\end{itemize}
		
	\end{columns}
	
\end{frame}


\begin{frame}{Een branch maken op GitHub}
	
	\begin{columns}[T]
		
		\column{0.48\textwidth}
		
		\textcolor{VSCodeBlue}{\textbf{Wat doen we?}}
		
		\begin{itemize}
			\item Maak branch \texttt{vakken} op GitHub.
			\item Maak bestand \texttt{vakken.txt}.
			\item Voeg inhoud toe.
			\item Commit rechtstreeks op GitHub.
		\end{itemize}
		
		\column{0.48\textwidth}
		
		\textcolor{GitOrange}{\textbf{Waarom?}}
		
		Niet alle branches ontstaan lokaal.
		
		\medskip
		
		\textcolor{GitHubGray}{\textbf{Resultaat}}
		
		\begin{itemize}
			\item branch \texttt{vakken} bestaat enkel online
			\item lokaal kent Git deze branch nog niet
			\item repositories bevatten verschillende informatie
		\end{itemize}
		
	\end{columns}
	
\end{frame}


\begin{frame}{Fetch -- online informatie ophalen}
	
	\begin{columns}[T]
		
		\column{0.48\textwidth}
		
		\textcolor{VSCodeBlue}{\textbf{Wat doen we?}}
		
		\begin{itemize}
			\item Bekijk de gekende branches.
			\item Klik \texttt{Fetch From All Remotes}.
			\item Bekijk opnieuw de branchlijst.
		\end{itemize}
		
		\column{0.48\textwidth}
		
		\textcolor{GitOrange}{\textbf{Waarom?}}
		
		Git controleert welke branches en commits online bestaan.
		
		\medskip
		
		\textcolor{GitHubGray}{\textbf{Resultaat}}
		
		\begin{itemize}
			\item \texttt{origin/vakken} wordt zichtbaar
			\item geen bestanden gewijzigd
			\item geen commits opgehaald
		\end{itemize}
		
	\end{columns}
	
\end{frame}


\begin{frame}{Pull -- een branch ophalen}
	
	\begin{columns}[T]
		
		\column{0.48\textwidth}
		
		\textcolor{VSCodeBlue}{\textbf{Wat doen we?}}
		
		\begin{itemize}
			\item Kies \texttt{Create New Local Branch From}.
			\item Selecteer \texttt{origin/vakken}.
			\item Maak lokale branch \texttt{vakken}.
			\item Schakel naar deze branch.
		\end{itemize}
		
		\column{0.48\textwidth}
		
		\textcolor{GitOrange}{\textbf{Waarom?}}
		
		De branch bestond tot nu toe uitsluitend online.
		
		\medskip
		
		\textcolor{GitHubGray}{\textbf{Resultaat}}
		
		\begin{itemize}
			\item lokale branch \texttt{vakken}
			\item bestand \texttt{vakken.txt} verschijnt
			\item branch kan lokaal gebruikt worden
		\end{itemize}
		
	\end{columns}
	
\end{frame}


\begin{frame}{Push -- branch publiceren}
	
	\begin{columns}[T]
		
		\column{0.48\textwidth}
		
		\textcolor{VSCodeBlue}{\textbf{Wat doen we?}}
		
		\begin{itemize}
			\item Schakel naar \texttt{woonplaats}.
			\item Kies \texttt{Publish Branch}.
			\item Controleer GitHub.
		\end{itemize}
		
		\column{0.48\textwidth}
		
		\textcolor{GitOrange}{\textbf{Waarom?}}
		
		Een lokale branch kan gedeeld worden via GitHub.
		
		\medskip
		
		\textcolor{GitHubGray}{\textbf{Resultaat}}
		
		\begin{itemize}
			\item drie branches zichtbaar op GitHub
			\item lokale en remote branch gesynchroniseerd
			\item cloudlabels verschijnen in Graph
		\end{itemize}
		
	\end{columns}
	
\end{frame}

\begin{frame}{Terminologie}
	
	\begin{columns}[T]
		
		\column{0.48\textwidth}
		
		\textcolor{VSCodeBlue}{\textbf{Nieuwe begrippen}}
		
		\begin{itemize}
			\item Branch
			\item Main
			\item Fetch
			\item Pull
			\item Push
		\end{itemize}
		
		\column{0.48\textwidth}
		
		\textcolor{GitHubGray}{\textbf{Samenvatting}}
		
		\begin{itemize}
			\item Branch = alternatieve versie
			\item Main = stabiele hoofdbranch
			\item Fetch = online informatie bekijken
			\item Pull = wijzigingen ophalen
			\item Push = wijzigingen publiceren
		\end{itemize}
		
		\medskip
		
		\textcolor{GitOrange}{\textbf{Volgende stap}}
		
		Branches opnieuw samenvoegen via merge.
		
	\end{columns}
	
\end{frame}

\section{Mergen van branches}

\begin{frame}{Samenvoegen van branches}
	
	\begin{columns}[T]
		
		\column{0.48\textwidth}
		
		\textcolor{VSCodeBlue}{\textbf{Situatie}}
		
		Momenteel bestaan meerdere branches:
		
		\begin{itemize}
			\item \texttt{main}
			\item \texttt{woonplaats}
			\item \texttt{vakken}
		\end{itemize}
		
		\column{0.48\textwidth}
		
		\textcolor{GitOrange}{\textbf{Waarom mergen?}}
		
		Branches laten onafhankelijk werk toe.
		
		Uiteindelijk moeten de wijzigingen opnieuw worden samengebracht.
		
		\medskip
		
		\textcolor{GitHubGray}{\textbf{Soorten merges}}
		
		\begin{itemize}
			\item nieuwe bestanden
			\item verschillende wijzigingen
			\item conflicterende wijzigingen
		\end{itemize}
		
	\end{columns}
	
\end{frame}
%%%%%%%%%%%%%%%%%%%%%%%%%%%%%%%%%%%%%%%%%%%%%%%%%%

\begin{frame}{Lokale merge van nieuwe bestanden}
	
	\begin{columns}[T]
		
		\column{0.48\textwidth}
		
		\textcolor{VSCodeBlue}{\textbf{Wat doen we?}}
		
		\begin{itemize}
			\item Schakel naar \texttt{main}.
			\item Open Command Palette.
			\item Kies \texttt{Git: Merge}.
			\item Selecteer branch \texttt{woonplaats}.
			\item Push daarna naar GitHub.
		\end{itemize}
		
		\column{0.48\textwidth}
		
		\textcolor{GitOrange}{\textbf{Waarom?}}
		
		De branch bevat een bestand dat ontbreekt in \texttt{main}.
		
		\medskip
		
		\textcolor{GitHubGray}{\textbf{Resultaat}}
		
		\begin{itemize}
			\item \texttt{woonplaatsen.txt} verschijnt
			\item merge commit op \texttt{main}
			\item lokale merge zichtbaar in Graph
			\item na push ook zichtbaar op GitHub
		\end{itemize}
		
	\end{columns}
	
\end{frame}

Remote merge van nieuwe bestanden
\begin{frame}{Remote merge via Pull Request}
	
	\begin{columns}[T]
		
		\column{0.48\textwidth}
		
		\textcolor{VSCodeBlue}{\textbf{Wat doen we?}}
		
		\begin{itemize}
			\item Open Pull Requests.
			\item Maak een nieuwe Pull Request.
			\item Kies:
			\begin{itemize}
				\item base = \texttt{main}
				\item compare = \texttt{vakken}
			\end{itemize}
			\item Create Pull Request.
			\item Merge Pull Request.
			\item Fetch en Pull in VSCode.
		\end{itemize}
		
		\column{0.48\textwidth}
		
		\textcolor{GitOrange}{\textbf{Waarom?}}
		
		Een contributor stelt wijzigingen voor.
		
		Een reviewer beslist of deze worden overgenomen.
		
		\medskip
		
		\textcolor{GitHubGray}{\textbf{Resultaat}}
		
		\begin{itemize}
			\item branch \texttt{vakken} geïntegreerd
			\item merge eerst op GitHub
			\item lokale repository loopt achter
			\item Fetch + Pull synchroniseren opnieuw
		\end{itemize}
		
	\end{columns}
	
\end{frame}

\begin{frame}{Contributor en reviewer}
	
	\begin{columns}[T]
		
		\column{0.48\textwidth}
		
		\textcolor{VSCodeBlue}{\textbf{Contributor}}
		
		\begin{itemize}
			\item werkt op eigen branch
			\item maakt commits
			\item opent Pull Request
		\end{itemize}
		
		\column{0.48\textwidth}
		
		\textcolor{GitHubGray}{\textbf{Reviewer}}
		
		\begin{itemize}
			\item bekijkt wijzigingen
			\item controleert commits
			\item beslist over merge
		\end{itemize}
		
		\medskip
		
		\textcolor{GitOrange}{\textbf{Workflow}}
		
		\begin{center}
			
			Branch
			
			\(\downarrow\)
			
			Pull Request
			
			\(\downarrow\)
			
			Review
			
			\(\downarrow\)
			
			Merge
			
		\end{center}
		
	\end{columns}
	
\end{frame}


\begin{frame}{Merge zonder conflicten}
	
	\begin{columns}[T]
		
		\column{0.48\textwidth}
		
		\textcolor{VSCodeBlue}{\textbf{Wat doen we?}}
		
		\begin{itemize}
			\item Maak branch \texttt{programmeertaal}.
			\item Voeg favoriete taal toe.
			\item Commit.
			\item Voeg in \texttt{main} studentnummer toe.
			\item Commit.
			\item Merge beide branches.
		\end{itemize}
		
		\column{0.48\textwidth}
		
		\textcolor{GitOrange}{\textbf{Waarom lukt dit?}}
		
		De wijzigingen gebeuren op verschillende plaatsen.
		
		\medskip
		
		\textcolor{GitHubGray}{\textbf{Resultaat}}
		
		\begin{itemize}
			\item studentnummer aanwezig
			\item programmeertaal aanwezig
			\item Git combineert automatisch
			\item geen menselijke tussenkomst nodig
		\end{itemize}
		
	\end{columns}
	
\end{frame}


\begin{frame}{Automatisch combineren}
	
	\begin{columns}[T]
		
		\column{0.48\textwidth}
		
		\textcolor{VSCodeBlue}{\textbf{Situatie}}
		
		Branch A:
		
		\begin{itemize}
			\item nieuwe lijn bovenaan
		\end{itemize}
		
		Branch B:
		
		\begin{itemize}
			\item nieuwe lijn onderaan
		\end{itemize}
		
		\column{0.48\textwidth}
		
		\textcolor{GitOrange}{\textbf{Analyse}}
		
		Git vergelijkt:
		
		\begin{itemize}
			\item bronbranch
			\item doelbranch
			\item gemeenschappelijke voorouder
		\end{itemize}
		
		\medskip
		
		\textcolor{GitHubGray}{\textbf{Gevolg}}
		
		Git weet dat beide wijzigingen kunnen samen bestaan.
		
	\end{columns}
	
\end{frame}

\begin{frame}{Een merge met conflict}
	
	\begin{columns}[T]
		
		\column{0.48\textwidth}
		
		\textcolor{VSCodeBlue}{\textbf{Wat doen we?}}
		
		\begin{itemize}
			\item Maak branch \texttt{conflict}.
			\item Verander woonplaats naar
			\texttt{Minas Tirith}.
			\item Commit.
			\item Verander in \texttt{main}
			woonplaats naar
			\texttt{Camelot}.
			\item Commit.
			\item Merge.
		\end{itemize}
		
		\column{0.48\textwidth}
		
		\textcolor{GitOrange}{\textbf{Probleem}}
		
		Dezelfde lijn werd twee keer aangepast.
		
		\medskip
		
		\textcolor{GitHubGray}{\textbf{Resultaat}}
		
		Git weet niet welke versie juist is.
		
		Het mergeproces stopt.
		
	\end{columns}
	
\end{frame}


\begin{frame}{Conflict oplossen}
	
	\begin{columns}[T]
		
		\column{0.48\textwidth}
		
		\textcolor{VSCodeBlue}{\textbf{Wat doen we?}}
		
		\begin{itemize}
			\item Bekijk conflict in editor.
			\item Kies:
			\begin{itemize}
				\item Current
				\item Incoming
				\item Both
			\end{itemize}
			\item Stage wijziging.
			\item Commit.
			\item Push.
		\end{itemize}
		
		\column{0.48\textwidth}
		
		\textcolor{GitOrange}{\textbf{Waarom?}}
		
		Git kan de inhoud niet zelf kiezen.
		
		\medskip
		
		\textcolor{GitHubGray}{\textbf{Resultaat}}
		
		\begin{itemize}
			\item conflict opgelost
			\item merge voltooid
			\item repository opnieuw consistent
		\end{itemize}
		
	\end{columns}
	
\end{frame}


\begin{frame}{Terminologie}
	
	\begin{columns}[T]
		
		\column{0.48\textwidth}
		
		\textcolor{VSCodeBlue}{\textbf{Nieuwe begrippen}}
		
		\begin{itemize}
			\item Merge commit
			\item Bronbranch
			\item Doelbranch
			\item Pull Request
			\item Contributor
			\item Reviewer
			\item Merge conflict
		\end{itemize}
		
		\column{0.48\textwidth}
		
		\textcolor{GitHubGray}{\textbf{Samenvatting}}
		
		\begin{itemize}
			\item wijzigingen samenbrengen
			\item lokaal of via GitHub
			\item automatische merge indien mogelijk
			\item menselijke beslissing bij conflict
		\end{itemize}
		
		\medskip
		
		\textcolor{GitOrange}{\textbf{Volgende stap}}
		
		Werk organiseren met Issues en Assignees.
		
	\end{columns}
	
\end{frame}


% ======================================================================
\section{Werkverdeling: issues en assignees}
% ======================================================================

\begin{frame}{Issues en assignees}
	\begin{columns}[T]
		\column{0.48\textwidth}
		\textbf{Nieuwe mogelijkheden}
		
		\begin{itemize}
			\item Issues
			\item Feature requests
			\item Assignees
			\item Taakverdeling
		\end{itemize}
		
		\column{0.48\textwidth}
		\textbf{Waarom?}
		
		Bij grotere projecten werken verschillende personen tegelijk.
		
		\begin{itemize}
			\item Taken opsplitsen
			\item Verantwoordelijkheden vastleggen
			\item Voorkomen dat meerdere personen
			hetzelfde probleem oplossen
		\end{itemize}
	\end{columns}
\end{frame}

\begin{frame}{Een issue en branch aanmaken}
	\begin{columns}[T]
		\column{0.48\textwidth}
		\textbf{Studenten}
		
		\begin{itemize}
			\item Open \texttt{Issues}
			\item Klik \texttt{New Issue}
			\item Maak issue \texttt{1-film}
			\item Open het issue
			\item Klik \texttt{Create branch}
			\item Maak branch \texttt{1-film}
		\end{itemize}
		
		\column{0.48\textwidth}
		\textbf{Opzet en resultaat}
		
		Een taak wordt expliciet geregistreerd.
		
		\begin{itemize}
			\item Openstaand issue verschijnt
			\item Nieuwe branch gekoppeld aan issue
			\item Voorbereiding op taakuitvoering
			\item Branch bestaat voorlopig enkel op GitHub
		\end{itemize}
	\end{columns}
\end{frame}

\begin{frame}{Een assignee toekennen}
	\begin{columns}[T]
		\column{0.48\textwidth}
		\textbf{Studenten}
		
		\begin{itemize}
			\item Open het issue
			\item Klik op \texttt{Assignees}
			\item Selecteer één of meer contributors
		\end{itemize}
		
		\column{0.48\textwidth}
		\textbf{Opzet en resultaat}
		
		Verantwoordelijkheden worden toegewezen.
		
		\begin{itemize}
			\item Zichtbaar wie een taak uitvoert
			\item Filters per contributor mogelijk
			\item Overzicht van openstaande taken
		\end{itemize}
	\end{columns}
\end{frame}

\begin{frame}{De wijziging uitvoeren}
	\begin{columns}[T]
		\column{0.48\textwidth}
		\textbf{Studenten}
		
		\begin{itemize}
			\item Open branch \texttt{1-film}
			\item Pas \texttt{gegevens.txt} aan
			\item Voeg favoriete film toe
			\item Maak een commit
			\item Maak een pull request
		\end{itemize}
		
		\column{0.48\textwidth}
		\textbf{Opzet en resultaat}
		
		De contributor werkt het issue uit.
		
		\begin{itemize}
			\item Probleem opgelost in de issue-branch
			\item Hoofdbranch nog ongewijzigd
			\item Wijzigingen klaar voor review
		\end{itemize}
	\end{columns}
\end{frame}

\begin{frame}{Het issue sluiten}
	\begin{columns}[T]
		\column{0.48\textwidth}
		\textbf{Studenten}
		
		\begin{itemize}
			\item Reviewer behandelt de pull request
			\item Merge uitvoeren
			\item Eventueel issue heropenen
			\item Eventueel issue manueel sluiten
		\end{itemize}
		
		\column{0.48\textwidth}
		\textbf{Opzet en resultaat}
		
		Controleren of de taak echt opgelost is.
		
		\begin{itemize}
			\item Branch wordt gemerged
			\item Issue verhuist naar
			\texttt{Closed Issues}
			\item Historiek van oplossingen blijft bewaard
			\item Heropenen blijft mogelijk
		\end{itemize}
	\end{columns}
\end{frame}

\begin{frame}{Terminologie}
	\begin{columns}[T]
		\column{0.48\textwidth}
		\textbf{Nieuwe begrippen}
		
		\begin{itemize}
			\item Issue
			\item Bug
			\item Feature request
			\item Assignee
			\item Open issue
			\item Closed issue
		\end{itemize}
		
		\column{0.48\textwidth}
		\textbf{Betekenis}
		
		\begin{itemize}
			\item Issue = taak of probleem
			\item Bug = fout
			\item Feature request = nieuwe functie
			\item Assignee = verantwoordelijke
			\item Open = nog uit te voeren
			\item Closed = afgewerkt
		\end{itemize}
	\end{columns}
\end{frame}

\begin{frame}{Workflow met issues}
	\centering
	
	\Large
	
	Issue
	
	\vspace{0.4cm}
	
	\(\Downarrow\)
	
	\vspace{0.4cm}
	
	Branch
	
	\vspace{0.4cm}
	
	\(\Downarrow\)
	
	\vspace{0.4cm}
	
	Commit(s)
	
	\vspace{0.4cm}
	
	\(\Downarrow\)
	
	\vspace{0.4cm}
	
	Pull Request
	
	\vspace{0.4cm}
	
	\(\Downarrow\)
	
	\vspace{0.4cm}
	
	Merge
	
	\vspace{0.4cm}
	
	\(\Downarrow\)
	
	\vspace{0.4cm}
	
	Issue sluiten
	
\end{frame}

\section{De teamopdracht, repository en workflows}

%%%%%%%%%%%%%%%%%%%%%%%%%%%%%%%%%%%%%%%%%%%%%%%%%%
\begin{frame}{De teamopdracht: overzicht}
	
	\begin{columns}[T]
		
		\column{0.48\textwidth}
		
		\textcolor{VSCodeBlue}{\textbf{Wat gebeurt er?}}
		
		\begin{itemize}
			\item Iedereen werkt in dezelfde teamrepository.
			\item Taken worden verdeeld.
			\item Werk wordt via branches samengebracht.
			\item De teamleider bewaakt de hoofdbranch.
		\end{itemize}
		
		\column{0.48\textwidth}
		
		\textcolor{GitOrange}{\textbf{Waarom?}}
		
		Versiebeheer wordt pas echt nuttig wanneer meerdere personen samenwerken.
		
		\medskip
		
		\textcolor{GitHubGray}{\textbf{Doel}}
		
		\begin{itemize}
			\item geen bestanden verloren
			\item geen overschreven werk
			\item duidelijke taakverdeling
			\item controleerbare geschiedenis
		\end{itemize}
		
	\end{columns}
	
\end{frame}
%%%%%%%%%%%%%%%%%%%%%%%%%%%%%%%%%%%%%%%%%%%%%%%%%%


\begin{frame}{Team repository}
	
	\begin{columns}[T]
		
		\column{0.48\textwidth}
		
		\textcolor{VSCodeBlue}{\textbf{Wat doen we?}}
		
		\begin{itemize}
			\item Word lid van repository \texttt{mainX}.
			\item Clone de repository.
			\item Werk vanuit VSCode.
			\item Gebruik branches voor nieuwe bijdragen.
		\end{itemize}
		
		\column{0.48\textwidth}
		
		\textcolor{GitOrange}{\textbf{Waarom?}}
		
		De teamrepository vormt de centrale plaats waar alle resultaten samenkomen.
		
		\medskip
		
		\textcolor{GitHubGray}{\textbf{Resultaat}}
		
		\begin{itemize}
			\item alle teamleden werken op dezelfde basis
			\item gedeelde geschiedenis
			\item centrale opslag op GitHub
			\item teamleider beheert integratie
		\end{itemize}
		
	\end{columns}
	
\end{frame}

\begin{frame}{Vergaderverslagen}
	
	\begin{columns}[T]
		
		\column{0.48\textwidth}
		
		\textcolor{VSCodeBlue}{\textbf{Wat doen we?}}
		
		\begin{itemize}
			\item Werk in branch
			\texttt{0-vergaderingenX}.
			\item Verslaggever schrijft verslag.
			\item Dien een Pull Request in.
			\item Teamleider merge na goedkeuring.
		\end{itemize}
		
		\column{0.48\textwidth}
		
		\textcolor{GitOrange}{\textbf{Waarom?}}
		
		Een verslag is pas definitief wanneer het volledige team akkoord gaat.
		
		\medskip
		
		\textcolor{GitHubGray}{\textbf{Resultaat}}
		
		\begin{itemize}
			\item conceptversies gescheiden van hoofdbranch
			\item eenvoudige feedback
			\item duidelijke historiek van beslissingen
		\end{itemize}
		
	\end{columns}
	
\end{frame}

Wekelijkse branches
\begin{frame}{Wekelijkse branches}
	
	\begin{columns}[T]
		
		\column{0.48\textwidth}
		
		\textcolor{VSCodeBlue}{\textbf{Wat doen we?}}
		
		\begin{itemize}
			\item Werk telkens in de branch van de week.
			\item Voorbeelden:
			\begin{itemize}
				\item \texttt{oefeningenX}
				\item \texttt{implementatieX}
				\item \texttt{verslagX}
				\item \texttt{posterX}
			\end{itemize}
		\end{itemize}
		
		\column{0.48\textwidth}
		
		\textcolor{GitOrange}{\textbf{Waarom?}}
		
		Elke deeltaak krijgt een eigen ontwikkelomgeving.
		
		\medskip
		
		\textcolor{GitHubGray}{\textbf{Resultaat}}
		
		\begin{itemize}
			\item minder conflicten
			\item duidelijk overzicht
			\item deelresultaten kunnen afzonderlijk groeien
		\end{itemize}
		
	\end{columns}
	
\end{frame}

\begin{frame}{Issues voor alle deeltaken}
	
	\begin{columns}[T]
		
		\column{0.48\textwidth}
		
		\textcolor{VSCodeBlue}{\textbf{Wat doen we?}}
		
		\begin{itemize}
			\item Maak voor elke taak een issue.
			\item Ken een assignee toe.
			\item Werk aan één issue tegelijk.
			\item Gebruik sub-issues indien nodig.
		\end{itemize}
		
		\column{0.48\textwidth}
		
		\textcolor{GitOrange}{\textbf{Waarom?}}
		
		Issues maken zichtbaar:
		
		\begin{itemize}
			\item wat nog moet gebeuren;
			\item wie waaraan werkt;
			\item welke taken voltooid zijn.
		\end{itemize}
		
		\medskip
		
		\textcolor{GitHubGray}{\textbf{Resultaat}}
		
		\begin{itemize}
			\item duidelijke taakverdeling
			\item minder dubbel werk
			\item beter projectoverzicht
		\end{itemize}
		
	\end{columns}
	
\end{frame}

\begin{frame}{Commit, push en merge}
	
	\begin{columns}[T]
		
		\column{0.48\textwidth}
		
		\textcolor{VSCodeBlue}{\textbf{Wat doen we?}}
		
		\begin{itemize}
			\item Commit regelmatig.
			\item Push afgewerkt werk.
			\item Meld problemen via issues.
			\item Gebruik Pull Requests.
			\item Merge pas na controle.
		\end{itemize}
		
		\column{0.48\textwidth}
		
		\textcolor{GitOrange}{\textbf{Goede praktijk}}
		
		\begin{itemize}
			\item één deelprobleem tegelijk
			\item duidelijke commit messages
			\item niet uren lokaal werken
			\item tijdig synchroniseren
		\end{itemize}
		
		\medskip
		
		\textcolor{GitHubGray}{\textbf{Resultaat}}
		
		\begin{itemize}
			\item werk gaat niet verloren
			\item teamgenoten zien voortgang
			\item begeleiding kan opvolgen
			\item eindproduct groeit gecontroleerd
		\end{itemize}
		
	\end{columns}
	
\end{frame}


\end{document}